
\medskip
\noindent\SGARefTarget{sga3:viii:corollary:7:5}\textbf{Corollary 7.5.}
\textit{Suppose in addition that \(G_0\) is smooth over \(k\) and that
\(p\operatorname{id}_{G_0}\) is flat. Equivalently, by the structure theory
of algebraic groups over an algebraically closed field \(\overline{k}\),
\(G_{\overline{k}}\) contains no subgroup isomorphic to the additive group.
Then, for every \(\nu\geq\nu_0\) and every \(n>0\), the group
\({}_{p^\nu}H_n\) is flat over \(S_n\).}

Indeed, \(p^{\nu_0}\operatorname{id}_{Q_n}=0\) implies that
\(p^\nu\operatorname{id}_{H_n}\) factors through \(G_n\), giving the
commutative diagram
\SGARefTarget{sga3:viii:diagram:7:5}
\[
\begin{tikzcd}[column sep=huge,row sep=large]
H_n
  \arrow[r,"p^\nu\operatorname{id}_{H_n}"]
  \arrow[dr,"v_n"]
&
H_n
\\
G_n
  \arrow[u,"u_n"]
  \arrow[r,"p^\nu\operatorname{id}_{G_n}"]
&
G_n
  \arrow[u,"u_n"']
.
\end{tikzcd}
\]
We claim that \(v_n\) is flat. Since \(H_n\) and \(G_n\) are flat over
\(S_n\), it is enough to check that \(v_0\) is flat (SGA~1, IV, 5.9).
Thus we may take \(n=0\), so that \(S_n=\operatorname{Spec}(k)\). Since
\(p\operatorname{id}_{G_0}\), and hence
\(p^\nu\operatorname{id}_{G_0}\), is flat, its image is an induced open
subgroup \(G'_0\) of \(G_0\). The map \(u_0:G_0\to H_0\) is surjective, so
\(v_0\) takes its values in \(G'_0\) and may be regarded as a homomorphism
to \(G'_0\). Its composite with \(u_0\) is an epimorphism; hence \(v_0\)
itself is an epimorphism, and therefore a flat homomorphism into \(G'_0\),
and thus into \(G_0\). Consequently \(v_n\) is flat and
\(\ker v_n={}_{p^\nu}H_n\)%
\begingroup
\renewcommand{\thefootnote}{\(\dagger\)}
\footnote[0]{English editor's note: the source prints
\({}_{p^\nu}H\) over \(S\) here and \({}_{p^\nu}H\) over \(S_n\) in
the proof of \SGARefLink{sga3:viii:lemma:7:7}{Lemma 7.7}. The subscripts \(n\) are supplied in both places
because those arguments concern the base-changed family over \(S_n\).}
\endgroup
is flat over \(S_n\). \qed

\medskip
\noindent\SGARefTarget{sga3:viii:remark:7:5:after:01}\textbf{Remark.}
We did not explicitly use the assumption that \(G_0\) is smooth over \(k\).
It is, however, an easy consequence of the flatness of
\(p\operatorname{id}_{G_0}\)
\emph{[the source prints \(p\operatorname{id}_{G}\)]}; this is why the
condition was stated explicitly
in Corollary~\SGARefLink{sga3:viii:corollary:7:5}{7.5}.

\medskip
\noindent\SGARefTarget{sga3:viii:lemma:7:6}\textbf{Lemma 7.6.}
\textit{Let \(u:G\to H\) be a surjective monomorphism of commutative
algebraic groups over a field \(k\) of characteristic \(p>0\), and put
\(Q=H/G\), a purely infinitesimal group. There is an integer \(\nu_1\) such
that, for every \(\nu\geq\nu_1\), the sequence}
\[
  0\longrightarrow{}_{p^\nu}G
  \longrightarrow{}_{p^\nu}H
  \longrightarrow Q\longrightarrow0
\]
\textit{is exact.}

It is enough to establish exactness at \(Q\), which in turn amounts to the
injectivity of the homomorphism
\SGARefTarget{sga3:viii:eq:7:6:x}
\begin{equation}
  \tag{x}
  \mathscr O_{Q,e}\longrightarrow
  \mathscr O_{{}_{p^\nu}H,e}
\end{equation}
of local rings at the identity elements. Recall that \(Q\) is
set-theoretically supported at the single point \(e\). There is a natural
homomorphism
\SGARefTarget{sga3:viii:eq:7:6:xx}
\begin{equation}
  \tag{xx}
  \mathscr O_{{}_{p^\nu}H,e}
  \longrightarrow
  \mathscr O_{H,e}/\mathfrak m^{p^\nu},
\end{equation}
where \(\mathfrak m\) is the augmentation ideal, equivalently the maximal
ideal, of \(\mathscr O_{H,e}\). Indeed,
\(p^\nu\operatorname{id}_H\) vanishes on the kernel of the iterated
Frobenius homomorphism \(F^\nu\). The composite of
\SGARefLink{sga3:viii:eq:7:6:x}{(x)} and
\SGARefLink{sga3:viii:eq:7:6:xx}{(xx)} is the
natural composite
\SGARefTarget{sga3:viii:eq:7:6:xxx}
\begin{equation}
  \tag{xxx}
  \mathscr O_{Q,e}\longrightarrow
  \mathscr O_{H,e}\longrightarrow
  \mathscr O_{H,e}/\mathfrak m^{p^\nu}.
\end{equation}
Now \(\mathscr O_{Q,e}\to\mathscr O_{H,e}\) is injective, because
\(H\to Q\) is an epimorphism and therefore flat. Moreover,
\(\mathscr O_{Q,e}\) is Artinian, while
\(\bigcap_\nu\mathfrak m^{p^\nu}=0\) in \(\mathscr O_{H,e}\), so the
intersection of the corresponding ideals in \(\mathscr O_{Q,e}\) is also
zero. One of them is therefore zero. Thus
\SGARefLink{sga3:viii:eq:7:6:xxx}{(xxx)} is
injective, and a fortiori so is
\SGARefLink{sga3:viii:eq:7:6:x}{(x)}.

\medskip
\noindent\SGARefTarget{sga3:viii:lemma:7:7}\textbf{Lemma 7.7.}
\textit{Under the hypotheses of \SGARefLink{sga3:viii:corollary:7:5}{7.5}, there is an integer \(\nu_1\) such
that, for every \(\nu\geq\nu_1\) and every \(n\), the sequence of
\(S_n\)-groups}
\[
  0\longrightarrow{}_{p^\nu}G_n
  \longrightarrow{}_{p^\nu}H_n
  \xrightarrow{\,w_n\,}Q_n
  \longrightarrow0
\]
\textit{is exact. More precisely, \(w_n\) is faithfully flat and its kernel
is \({}_{p^\nu}G_n\).}

Take \(\nu_1\) as in Lemma~\SGARefLink{sga3:viii:lemma:7:6}{7.6}, applied
to \(u_0:G_0\to H_0\), and require \(\nu_1\geq\nu_0\), where
\(p^{\nu_0}=\operatorname{rank}Q_0\). It remains only to check that \(w_n\)
is faithfully flat. By Corollary~\SGARefLink{sga3:viii:corollary:7:5}{7.5},
\({}_{p^\nu}H_n\) is flat over \(S_n\), as is \(Q_n\); it is therefore
enough to check that
\[
  w_0:{}_{p^\nu}H_0\longrightarrow Q_0
\]
is faithfully flat, equivalently an epimorphism. This follows from the
choice of \(\nu_1\).

\medskip
\noindent\SGARefTarget{sga3:viii:corollary:7:8}\textbf{Corollary 7.8.}
\textit{Suppose in addition that \(H\) is separated over \(S\), or more
generally that \({}_{p^\nu}H\) is separated over \(S\) for every \(\nu\),
and that \({}_{p^\nu}G\) is finite over \(S\) for every \(\nu\). Then, for
\(\nu\geq\nu_1\), the group schemes \({}_{p^\nu}G\) and
\({}_{p^\nu}H\) are finite and flat over \(S\), and there is an exact
sequence}
\SGARefTarget{sga3:viii:corollary:7:8:display:exact-sequence:01}\[
  0\longrightarrow{}_{p^\nu}G
  \longrightarrow{}_{p^\nu}H
  \xrightarrow{\,w\,}Q
  \longrightarrow0.
\]

Because \(u:G\to H\) is surjective, the induced morphism
\({}_{p^\nu}G\to{}_{p^\nu}H\) is also surjective. The source is finite over
\(S\) and the target is separated over \(S\), hence the target is finite
over \(S\). To prove that \({}_{p^\nu}G\) and \({}_{p^\nu}H\) are flat over
\(S\), it is enough to prove this after every base change to \(S_n\), where
it is contained in Corollary~\SGARefLink{sga3:viii:corollary:7:5}{7.5}.
Finally, \SGARefLink{sga3:viii:corollary:7:8:display:exact-sequence:01}{the exact sequence above} is obtained from the exact sequences in
Lemma~\SGARefLink{sga3:viii:lemma:7:7}{7.7} as \(n\) varies.

Taking generic fibers in the sequence of
Corollary~\SGARefLink{sga3:viii:corollary:7:8}{7.8}, and recalling that
\(u_{s'}:G_{s'}\to H_{s'}\) is an isomorphism, gives that \(Q_{s'}\) is
the unit group. Since \(Q\) is flat over \(S\), \(Q\) itself is the unit
group; therefore so is \(Q_s\), and \(u_s:G_s\to H_s\) is an isomorphism.
We have proved:

\medskip
\noindent\SGARefTarget{sga3:viii:proposition:7:9}\textbf{Proposition 7.9.}
\textit{Let \(u:G\to H\) be a homomorphism of group preschemes of finite
presentation over a prescheme \(S\). Assume:}
\begin{enumerate}
  \item[(a)]\SGARefTarget{sga3:viii:proposition:7:9:item:a}
  \textit{\(u\) is a monomorphism;}
  \item[(b)]\SGARefTarget{sga3:viii:proposition:7:9:item:b}
  \textit{\(G\) is flat over \(S\);}
  \item[(c)]\SGARefTarget{sga3:viii:proposition:7:9:item:c}
  \textit{for every \(s\in S\) whose residue field has characteristic
  \(p>0\), let \(S'=\operatorname{Spec}(\mathscr O_{S,s})\), and let
  \(u':G'\to H'\) be obtained from \(u\) by base change to \(S'\). Then
  \(G'\) is commutative, the special fiber \(G'_0=G_s\) is smooth over
  \(\kappa(s)\), and, for every \(\nu>0\), the group
  \({}_{p^\nu}G'\) is finite over \(S'\) while \({}_{p^\nu}H'\) is
  separated over \(S'\).}
\end{enumerate}
\textit{Under these assumptions, \(u\) is an immersion.}

Indeed, the finiteness of \({}_pG'\) over \(S'\) implies the finiteness of
its special fiber over \(\kappa(s)\). Hence
\(G'_0\otimes_{\kappa(s)}\overline{\kappa(s)}\) contains no subgroup
isomorphic to the additive group, and the hypotheses of
Corollary~\SGARefLink{sga3:viii:corollary:7:5}{7.5} apply.

\medskip
\noindent\SGARefTarget{sga3:viii:remark:7:10}\textbf{Remark 7.10.}
Examples of M. Raynaud (XVI, 1.1(a),\SGARefLink{sga3:viii:proposition:7:9:item:b}{(b)}) show that in condition \SGARefLink{sga3:viii:proposition:7:9:item:c}{(c)} one
cannot omit either the finiteness of the \({}_{p^\nu}G'\) over \(S'\) or
the separatedness of the \({}_{p^\nu}H'\) over \(S'\).

We now seek conditions under which \(u\) is a closed immersion. Retain the
hypotheses preceding Proposition~\SGARefLink{sga3:viii:proposition:7:9}{7.9},
but no longer suppose that \(u\) is surjective; assume only that it is an
isomorphism on generic fibers. We already know that \(u:G\to H\) is an open
immersion. The same is therefore true of \(u_0:G_0\to H_0\), whose image
contains the identity component \((H_0)^0\). Since \(G_0\) is smooth over
\(k\) and has no additive component after extending to \(\overline{k}\),
the same is true of \(H_0\). We use:

\medskip
\noindent\SGARefTarget{sga3:viii:lemma:7:11}\textbf{Lemma 7.11.}
\textit{Let \(H\) be a commutative algebraic group over a field \(k\), and
suppose that \(H\otimes_k\overline{k}\) contains no subgroup isomorphic to
\(\mathbf G_a\). If \(n=\deg(H/H^0)\), then}
\[
  {}_nH\longrightarrow H/H^0
\]
\textit{is surjective.}

We may assume that \(k\) is algebraically closed, in which case this follows
from the familiar fact that \(H^0(k)\) is a divisible group.

Return to \(u:G\to H\), and suppose that the \({}_nG\), \(n>0\), are proper
over \(S\), while the \({}_nH\) are separated over \(S\). The \({}_nH\)
are also flat over \(S\). Indeed, it is enough to check flatness at the
points above \(s\), and hence to check that
\(n\operatorname{id}_{H_0}\) is flat. As noted above, this is equivalent to
\((H_0)^0\) being smooth\footnote{Editor's note: the terminology has been
updated, replacing ``simple'' by ``smooth''; see, for example, A.
Grothendieck's note (1) in SGA~1, Expos\'e II.} over \(k\) and having no
\(\mathbf G_a\)-component over \(\overline{k}\). Since
\((H_0)^0=(G_0)^0\), this follows from the corresponding hypothesis on
\(G_0\).

Because \(H\) is separated over \(S\), so are the \({}_nH\). The morphism
\({}_nG\to{}_nH\) is proper and therefore has closed image. Its image
contains the generic fiber of \({}_nH\), because
\(u_{s'}:G_{s'}\to H_{s'}\) is an isomorphism. Since \({}_nH\) is flat
over \(S\), it is the schematic closure of its generic fiber. Hence
\({}_nG\to{}_nH\), and therefore \({}_nG_0\to{}_nH_0\), is surjective.
Together with \(G_0\supset(H_0)^0\), Lemma
\SGARefLink{sga3:viii:lemma:7:11}{7.11} now implies that
\(G_0\to H_0\) is surjective. Hence \(u\) is surjective. We obtain:

\medskip
\noindent\SGARefTarget{sga3:viii:proposition:7:12}\textbf{Proposition 7.12.}
\textit{With the notation of \SGARefLink{sga3:viii:proposition:7:9}{7.9}, suppose that conditions \SGARefLink{sga3:viii:proposition:7:9:item:a}{(a)} and \SGARefLink{sga3:viii:proposition:7:9:item:b}{(b)} hold,
together with the stronger condition (c\(^{\prime}\)) obtained from \SGARefLink{sga3:viii:proposition:7:9:item:c}{(c)} by
requiring, for every integer \(n>0\), not only those of the form \(p^\nu\),
that \({}_nG'\) be finite over \(S'\) and \({}_nH'\) separated over \(S'\).
Then \(u\) is a closed immersion.}

\medskip
\noindent\SGARefTarget{sga3:viii:remark:7:13}\textbf{Remark 7.13.}
\begin{enumerate}
  \item[(a)] The separatedness hypothesis on the \({}_nH\) easily implies
  that \(H\) itself is separated over \(S\). This is also formally contained
  in Proposition~\SGARefLink{sga3:viii:proposition:7:12}{7.12}, by taking
  \(G\) to be the unit \(S\)-group. If \(S\) is locally Noetherian, in
  Proposition~\SGARefLink{sga3:viii:proposition:7:9}{7.9} it is enough to
  assume that \(H\) is locally of finite type over \(S\), rather than of
  finite type; the proof applies verbatim. For
  Proposition~\SGARefLink{sga3:viii:proposition:7:12}{7.12}, one also
  assumes that the fibers of \(H\) are of finite type.

  \item[(b)]\SGAInlineRefTarget{sga3:viii:remark:7-13:item:b} Using Proposition~\SGARefLink{sga3:viii:proposition:7:12}{7.12},
  it is not difficult to prove that if \(u:G\to H\) is a monomorphism of
  group preschemes of finite presentation over \(S\), \(G\) is
  diagonalizable---or, more generally, of multiplicative type---and \(H\)
  is separated over \(S\), then \(u\) is a closed immersion. This is a
  satisfactory generalization of the first conclusion in \SGARefLink{sga3:viii:corollary:5:7}{5.7}. If \(G\) is
  smooth over \(S\), apply \SGARefLink{sga3:viii:proposition:7:12}{Proposition~7.12} directly; the general case
  reduces readily to this one. If \(H\) is not assumed separated over \(S\),
  one can still show that \(u\) is an immersion; for a torus \(G\), this
  also follows from the material that comes next.

  \item[(c)] If \(G\) is assumed to have connected fibers in
  Proposition~\SGARefLink{sga3:viii:proposition:7:9}{7.9}, the
  separatedness of \({}_{p^\nu}H'\) over \(S'\) may be omitted from
  condition \SGARefLink{sga3:viii:proposition:7:9:item:c}{(c)}. Indeed, after the reductions in the proof, \(H\) may be
  assumed flat over \(S\) and \(u\) bijective. Thus \(H\) has connected
  fibers, and Raynaud's theorem (VI\(_B\), 5.5) implies that \(H\) is
  separated over \(S\).\footnote{Editor's note: the attribution, which
  appeared in SGAD 1965, should be made explicit in Expos\'e VI\(_B\), §5,
  together with the changes between SGAD and Lecture Notes 151.}
\end{enumerate}

\SGARefTarget{sga3:viii:bibliography}\section*{Bibliography}
\addcontentsline{toc}{section}{Bibliography}
\footnote{Editor's note: additional references cited in this expos\'e.}

\begin{description}
  \item[\SGARefTarget{sga3:viii:bib:rg71}{[RG71]}]
  M. Raynaud and L. Gruson, \textit{Crit\`eres de platitude et de
  projectivit\'e}, Invent. Math. \textbf{13} (1971), 1--89.

  \item[\SGARefTarget{sga3:viii:bib:to70}{[TO70]}]
  J. Tate and F. Oort, \textit{Group schemes of prime order}, Ann. Sci.
  \'Ec. Norm. Sup\'er. (4) \textbf{3} (1970), 1--21.
\end{description}

